\section{Related literature}\label{sec:lit}
% =============================================================================================
% Source: docs/literature.md (entries, "Closest competitors", "Positioning for a CREST
% audience", "Citation map"). Only keys of references.bib (Sections 9 and 10 of the .bib hold
% the keys proposed in OUTLINE.md Section 13 and the six keys added in the referee revision:
% im_2018_quantitatively, gulrajani_2019_towards, sutherland_2017_generative,
% jitkrittum_2016_interpretable, meehan_2020_three, lucic_2018_gans; each checked against the
% publisher or arXiv record at integration, 2026-11-16). Positioning only: no
% results, no numbers. Never "first". Subsection labels: OUTLINE.md Section 12, plus
% sec:lit-summary (added here).
% =============================================================================================

This section positions the paper in five literatures; \cref{sec:lit-summary} summarises what is
and is not ours.

\subsection{Generative models for financial time series and their evaluation}
\label{sec:lit-synthetic}

Synthetic return data is used for back-testing, data sharing and hedging
\citep{assefa2020generating,kondratyev2019market,buehler2019deep,bailey2017probability}.
Generators include GANs \citep{wiese2020quant,takahashi2019modeling}, restricted Boltzmann
machines \citep{kondratyev2019market}, signature-based variational autoencoders
\citep{buehler2020data}, generators trained with signature losses
\citep{ni2021sig,liao2024sig,issa2023nonadversarial}, generic time-series GANs
\citep{esteban2017real,yoon2019time} and generators with use-specific losses such as portfolio
tail risk \citep{cont2026tail}. Parametric models remain natural benchmarks (GARCH(1,1), for
instance, is hard to beat out of sample; \citealp{hansen2005forecast}); together with a stationary
bootstrap they are the generators of this paper (\cref{sec:data-generators}).

Evaluation practice is heterogeneous \citep{stenger2024evaluation}. Most papers compare stylized
facts \citep{cont2001empirical} visually or through summary statistics
\citep{wiese2020quant,takahashi2019modeling,vyetrenko2020get}. Others report fidelity scores
\citep{alaa2022faithful} or discriminator scores: the held-out error of a classifier trained after
the fact \citep{yoon2019time}, a real-versus-simulated classification competition
\citep{francopedroso2018ets}, or a discriminator output used to calibrate a market simulator
\citep{storchan2021learning}. In machine learning, critics for several GAN objectives, the
Wasserstein one included, have been trained after the fact on training samples and the resulting
divergences evaluated on held-out samples as metrics for image generators
\citep{im_2018_quantitatively}, and \citet{gulrajani_2019_towards} define such neural-network
divergences on held-out data so that a generator that memorises its training set does not score
well. Such scores have no null distribution, and a low discriminative score
cannot be told apart from a weak classifier unless the power of the procedure is known; hence the
pre-specified power check in our decision rule (\cref{sec:method-decision}).

Formal testing of market generators is not new. \Citet{buehler2020data} argue for two-sample tests
as evaluation metrics and apply marginal Kolmogorov--Smirnov tests and the signature-MMD test of
\citet{chevyrev2022signature}, whose distribution-free threshold is valid for every null law and
hence conservative; they report no size or power analysis. \Citet{bai2026calibrating} come closer
to this paper: to calibrate non-identifiable simulators, among them a limit-order-book market
simulator, they compare real and simulated outputs through a feature map, possibly the output of a
trained WGAN critic, with Bonferroni-corrected Kolmogorov--Smirnov or mean-difference statistics
and asymptotic guarantees for independent outputs. A WGAN critic has thus already been used, with
frequentist guarantees, to validate a financial simulator against real data. Our goal differs (a
verdict on one fitted generator, not a set of eligible parameters), and so does the statistic:
their critic is a fixed pre-trained feature map (the most conservative of their features), ours is
trained on one half of the two samples and its mean difference on the other half is the statistic.
Dependent windows, a relative null and the likelihood-ratio oracle are outside their scope.

\subsection{Two-sample tests with learned test functions}\label{sec:lit-tests}

An integral probability metric (IPM) is $\sup_{f\in\mathcal F}\abs{Pf-Qf}$ over a class
$\mathcal F$ of test functions \citep{muller1997integral}; over the 1-Lipschitz functions it is
$\W$, by Kantorovich--Rubinstein duality \citep{villani2009optimal}. The critic of a Wasserstein GAN maximises the empirical
objective over a class of networks \citep{arjovsky2017wasserstein}, with a gradient penalty in
place of the Lipschitz constraint in WGAN-GP \citep{gulrajani2017improved}. Our critic, trained
alone, is a neural-network test function chosen to maximise a mean difference under a Lipschitz
penalty. An IPM over a restricted class is estimable at the parametric rate
\citep{sriperumbudur2012empirical}, whereas the empirical measure converges in $\W$ at rate
$n^{-1/d}$ for $d\ge3$ \citep{fournier2015rate}. A finite critic measures a weaker ``neural net distance''
\citep{arora2017generalization}, and with a soft penalty a good critic need not estimate $\W$ well
\citep{mallasto2019well}. Our statistic is therefore an IPM over trained critics, whose
conditional mean is bounded by the critic's Lipschitz constant times $\W$ between the standardised
window laws (\cref{prop:unbiased,rem:gp-soft}), not an estimate of $\W$. A trained critic is
also a random function of its seed and hyper-parameters, and large-scale comparisons of GANs
depend on random restarts and on the hyper-parameter search budget \citep{lucic_2018_gans}. We
therefore combine p-values over seeds by a rule that preserves the validity of the individual
p-values under arbitrary dependence between them (\cref{rem:combine-seeds}) and report how the
verdicts change with the critic's hyper-parameters (\cref{sec:robust}).

The MMD test compares kernel mean embeddings \citep{gretton2012kernel}.
\Citet{sutherland_2017_generative} choose the kernel on a training split by maximising an estimate
of a power proxy, test on the held-out split, and use the optimised MMD and its witness function
to criticise generative models; \citet{liu2020learning} extend this to deep kernels (MMD-D), and
\citet{jitkrittum_2016_interpretable} choose test locations by maximising a lower bound on power,
so that the selected features show where the two laws differ. These tests, like the evaluation
metrics of \cref{sec:lit-synthetic}, treat the samples as independent. The classifier
two-sample test (C2ST) of \citet{lopezpaz2017revisiting} tests whether held-out accuracy differs
from one half. Accuracy discards magnitude: in the Gaussian mean-shift problem accuracy tests are
only $1/\sqrt{\pi}$ as efficient as Hotelling-type tests \citep{kim2021classification}. The
held-out difference of mean logits \citep{cheng2022classification} keeps it; since the Bayes logit
is the log density ratio
\citep{goodfellow2014generative,nguyen2010estimating,cranmer2015approximating}, it is a learned
log-likelihood-ratio statistic, our \LRC{} competitor (\cref{sec:method-competitors}).

The \WC{} test instantiates a known template: learn a real-valued witness on one split and test the
equality of its means on the other, with a permutation null \citep{kubler2022witness}; see also
the location tests on learned representations of \citet{kirchler2020two}, the logit instance of
\citet{cheng2022classification} and the witness-based model criticism of
\citet{lloyd2015statistical}. The split-then-permute structure of a Wasserstein test is not ours
either: \citet{wang2022two} (linear version: \citealp{wang2021two}) train a kernel projector on one
split to maximise the Wasserstein distance between projected samples and permute on the held-out
split with the projector fixed. \Citet{hu2026neural} use a deep, approximately 1-Lipschitz
witness, as we do, but with learned projections, max-type aggregation and asymptotically pivotal
calibration for \iid{} data; see also \citet{ramdas2017wasserstein}. Our test is the template
with a WGAN-GP critic: the statistic is the held-out mean difference of one network (no transport
problem is solved on the evaluation split), calibrated by permutation for independent windows and
by a Driscoll--Kraay variance for dependent ones. We claim this instance, its validation and the
protocol around it, not the template. We do not show that it is more powerful than MMD-D, which is
not among our baselines (\cref{sec:concl-limitations}), and it does not beat the baselines we run
against every deviation: against changes of the one-dimensional marginal of returns, the pooled
Kolmogorov--Smirnov test is far more powerful than every window test in our simulations, although
it is invalid when returns within a window are dependent and blind to deviations that leave the
marginal unchanged (\cref{sec:sim-power}).

\Citet{chevyrev2022signature} show that (robust) signature moments characterise laws of
stochastic processes and build an MMD test on them
\citep[see also][]{lyons1998differential,kiraly2019kernels}. Signature discrepancies also serve as
training losses for market generators \citep{ni2021sig,liao2024sig,issa2023nonadversarial}, one
reason to evaluate with a test that was not the training objective. Our \SigMMD{} baseline is a
truncated (depth-4) signature-feature MMD in their spirit, without tensor normalisation; it does
not inherit their characteristic property.

A relative test compares two candidate models with one reference sample.
\Citet{bounliphone2016test} use the difference of two MMD estimates and their joint asymptotic
normality, which requires independent \iid{} samples and two candidates that both differ from the
reference, the MMD U-statistic being degenerate at equality; our MMD relative test follows them
and inherits this condition. See also \citet{jitkrittum2018informative} and
\citet{gerber2024likelihood} on deciding which of two simulators produced the data. Closest in
aim, \citet{zhou_2025_anchor} ask which of two distributions is closer to an anchor distribution:
they maximise the difference between the two anchored discrepancies over deep kernels, which also
selects the direction of the hypothesis, and then calibrate the selected test, for independent
samples. Our relative test differs in four respects. One of the two candidates is the real past,
so the question is whether a generator does better than historical simulation; the windows are
serially and cross-sectionally dependent; the discrepancy is the held-out mean difference of one
critic, whose witness can be inspected; and the test enters a decision rule with a pre-specified
power check.
% zhou_2025_anchor: description checked against the arXiv:2510.10477 abstract (anchor-based
% maximum discrepancy over a deep-kernel space; Phase I infers the hypothesis, Phase II tests it).
In our relative null, then, $\Rpast$ is one of the candidates: is $G$ as good a proxy for the future
$\Rnew$ as history is? With learned critics, the difference of two separately trained
discrepancies compares IPMs over two data-dependent classes, and at the boundary $\QG=\Ppast$
its size is at least the nominal level and can exceed it (\cref{prop:boundary-size}). The
anchored test trains one critic to separate $\Rnew$ from a balanced pool of synthetic and past
windows and compares its means on held-out synthetic and past windows; its permutation p-value is
exact at the boundary for every critic (\cref{prop:anchored}). The argument is the elementary
fixed-witness one; its value is exactness at the boundary, where the other two versions are not
exact. Exactness needs the synthetic and past windows to be exchangeable at the boundary
(\cref{ass:anc}). Serially dependent windows are not, and the Driscoll--Kraay version of the test
used on such panels has no such guarantee: its size is measured
(\cref{sec:sim-dependent,sec:emp-placebo}). A generator that copies its training data violates the
assumption as well. Data copying is a known failure of generator evaluation:
\citet{meehan_2020_three} test for it by comparing the distances of generated and of held-out real
samples to the training sample. Our stationary bootstrap copies blocks of $\Rpast$ by
construction, so we do not test for copying; we withhold its verdict (\INVALID{};
\cref{sec:method-relative,tab:e2-relative-boundary}).

\subsection{Specification testing and density forecast evaluation in econometrics}
\label{sec:lit-econometrics}

For an econometrician the natural frame is the specification testing of simulation-based models.
Models without a closed-form likelihood are estimated by simulated moments
\citep{mcfadden1989method,duffie1993simulated}, indirect inference
\citep{gourieroux1993indirect,smith1993estimating}, the efficient method of moments
\citep{gallant1996which} or simulated maximum likelihood \citep{fermanian2004nonparametric}; see
\citet{gourieroux1997simulation}, and \citet{bernton2019approximate} and \citet{cherief2022finite}
for Wasserstein ABC and minimum-MMD estimation under misspecification. Discriminators are accepted econometric objects: in the
adversarial estimator of \citet{kaji2023adversarial} a logistic discriminator gives optimally
weighted simulated moments and the oracle discriminator maximum likelihood \citep[see
also][]{kaji2023metropolis}. Our paper is a testing counterpart: we test a calibrated generator
instead of estimating one, use an IPM instead of cross-entropy, and work with serially dependent
windows, which their theory excludes. Their bridge from discriminators to maximum likelihood is the
estimation counterpart of our comparison of the \LRC{} statistic with the Neyman--Pearson oracle
(\cref{prop:bayes-logit,prop:oracle-bound}).

Specification tests of financial models compare the marginal density of a diffusion with a kernel
estimate \citep{aitsahalia1996testing} or test transition densities through probability integral
transforms \citep{hong2005nonparametric}, which requires the transition density. Closest to us,
\citet{corradi2005bootstrap} and \citet{bhardwaj2008simulation} compare the empirical (marginal,
respectively conditional) distribution of the data with that of data simulated from the fitted
model, covering stochastic volatility without a closed-form likelihood, with block-bootstrap
critical values that account for serial dependence and parameter estimation error. Our test
differs in three respects. (a) \emph{Null}: they test the model class at the pseudo-true
parameter, so estimation error must be corrected; we test the fitted generator (parameters
estimated on \TRAIN{} and held fixed, null $\QG=\Pnew$ on \TEST{}), so no correction is needed,
but a rejection does not separate misspecification from estimation error or parameter drift, which
the relative test is designed to absorb. (b) \emph{Statistic}: a fixed Kolmogorov-type functional
versus a learned function. (c) \emph{Calibration}: a block bootstrap versus permutation or a
Driscoll--Kraay variance, conditional on the fitted parameters and the critic.

In this language the critic is a test function chosen from the data. Consistent specification
tests use a whole class of test functions, such as exponential weights
\citep{bierens1990consistent} or neural hidden units indexed by a nuisance parameter present only
under the alternative \citep{stinchcombe1998consistent,lee1993testing}, with a supremum or integral
over the class and asymptotic calibration. The \WC{} test is a learned, sample-split relative: it
selects one function on a training split and calibrates on the other. The analogy concerns the test
functions only (their null is a conditional moment restriction), and the \WC{} test inherits the
power of the learned function, not the consistency of Bierens-type tests; it is consistent only
under conditions on training (\cref{prop:consistency,rem:c-realism}).

When a likelihood exists, density forecast evaluation supplies trusted benchmarks: probability
integral transforms are \iid{} uniform under a correct predictive density
\citep{diebold1998evaluating}, \citet{berkowitz2001testing} tests their normal transforms against a
Gaussian AR(1), the log score is strictly proper \citep{gneiting2007strictly}, and density forecasts
are compared by out-of-sample log-likelihood ratios with a HAC variance
\citep{amisano2007comparing}, in the spirit of \citet{vuong1989likelihood}. We use them for GBM,
Merton and GARCH-$t$ (\cref{sec:emp-likelihood}); Heston has no closed-form likelihood
\citep{aitsahalia2007maximum} and the stationary bootstrap no parametric likelihood. Likelihoods rank models by
Kullback--Leibler divergence and critics by an IPM, so disagreement is informative
(\cref{rem:kl-vs-w1}). The relative test is a loss-differential comparison
\citep{diebold1995comparing} that, like \citet{giacomini2006tests}, evaluates a forecasting method
(``use the real past''); its closest econometric relative is the simulation-based comparison of
misspecified diffusion models by \citet{corradi2011predictive}.

\subsection{Inference under dependence}\label{sec:lit-dependence}

A permutation test is exact when the pooled observations are exchangeable under the null
\citep{lehmann2022testing}, also with random permutations and the $+1$ correction
\citep{hemerik2018exact,phipson2010permutation}; with a critic fixed by the training half this
gives exact level for independent windows (\cref{prop:exact-level}). Studentised permutation tests
are asymptotically valid for the weaker null of equal means, for independent samples
\citep{janssen1997studentized,chung2013exact}; under time-series dependence the closest results
concern autocorrelation tests \citep{romano2022permutation}. Real return windows are not
exchangeable: consecutive windows share volatility regimes, and windows of different stocks on the
same dates share market shocks. We studentise the held-out mean difference with a Bartlett-kernel
HAC variance \citep{newey1987simple} and the bandwidth of \citet{andrews1991heteroskedasticity},
after summing centred critic values by date \citep{driscoll1998consistent}. The normal
approximation rests on a central limit theorem for mixing sequences \citep{herrndorf1984functional},
and GARCH and a class of stochastic-volatility models are $\beta$-mixing
\citep{carrasco2002mixing}. With few time periods HAC tests over-reject; fixed-$b$ critical values
\citep{kiefer_2005_new} are one correction, reported for comparison (\cref{rem:small-T}). The MMD
statistic is a degenerate U-statistic under the null rather than a mean, and
\citet{chwialkowski2014wild} calibrate it under dependence by a wild bootstrap with dependent
multipliers, which we use for MMD and \SigMMD{} (\cref{rem:wild}); block bootstraps
\citep{kunsch1989jackknife,politis1994stationary,lahiri2003resampling} are the alternative.

Two features are less standard. First, the critic is trained on data that depend on the evaluation
data. A ticker split puts both halves on the same dates and biases the statistic
(\cref{prop:ticker-bias}); a split by contiguous time blocks makes training and evaluation dates
disjoint, which removes the common-date bias but does not make the two halves independent, and
a persistent factor can still bias the statistic across adjacent blocks (\cref{cor:block-bias}). Under
mixing, \cref{lem:decoupling} bounds the remaining dependence across block boundaries only when a
gap separates the two halves; the implemented split has none, and under persistent volatility
adjacent blocks remain strongly dependent, so the effect of this dependence on size is measured,
not bounded (\cref{sec:sim-dependent,sec:emp-placebo}). Second, per-ticker tests are
dependent through the market factor: Benjamini--Hochberg control relies on positive regression
dependence, and arbitrary dependence needs the Benjamini--Yekutieli correction
\citep{benjamini1995controlling,benjamini2001control}.

\subsection{Portfolio choice and use-based evaluation}\label{sec:lit-portfolio}

Mean-variance portfolios \citep{markowitz1952portfolio} are very sensitive to estimation error
\citep{michaud1989markowitz,chopra1993effect,kan2007optimal}; shrinkage \citep{ledoit2004well},
long-only constraints \citep{jagannathan2003risk} and the equally weighted portfolio
\citep{demiguel2009optimal} are standard responses. A generator fitted on \TRAIN{} is a randomised
function of \TRAIN, so its output carries no more information about the future than \TRAIN{}
itself. \Citet{cetingoz2026synthetic} make this precise for Markowitz and minimum-variance
portfolios: the initial sample size limits what synthetic data can add, generating far more data
than was available can bias the estimated statistic, and generic generators fit high-variance
principal components that matter little for minimum-variance portfolios; they argue for
use-targeted evaluation. Synthetic-trained portfolios can thus beat a specific real-data rule, such
as the sample-covariance GMV portfolio, only when a parsimonious generator acts as a structured,
shrinkage-like estimator (\cref{prop:no-free-lunch}).

Judging synthetic data by a downstream decision is the train-on-synthetic, test-on-real (TSTR)
principle \citep{esteban2017real,yoon2019time}; in econometrics, \citet{athey2024using} judge
generated data by the ranking of estimators it produces, and evaluation criteria for generative
models are known to disagree \citep{theis2016note}. Our Markowitz study uses TSTR as the utility
measure but asks whether a distributional test verdict predicts the TSTR utility gap across
universes and generators, which \citet{cetingoz2026synthetic} do not address. Theory suggests at
most a weak link (\cref{sec:mk-theory}): GMV utility depends on second moments only, so critic
distance and utility can disagree in both directions, and under non-stationarity a generator that
reproduces the past better is not necessarily more useful (\cref{tab:e4-correlation}).

\subsection{Summary of the positioning}\label{sec:lit-summary}

None of the following is ours: a critic or neural-network divergence trained after the fact to
evaluate a generator \citep{im_2018_quantitatively,gulrajani_2019_towards}; a kernel learned on one
split to criticise a generative model \citep{sutherland_2017_generative}; a witness learned on one
split and tested on the other
\citep{kubler2022witness,kirchler2020two,cheng2022classification}; a split-then-permute
Wasserstein test \citep{wang2022two}; a WGAN critic used to validate a financial simulator
\citep{bai2026calibrating}; a formal two-sample test of a market generator
\citep{buehler2020data}; a specification test comparing real data with data simulated from the
fitted model, with dependence handled by a block bootstrap \citep{corradi2005bootstrap}. What the paper
adds is their combination and its validation, as listed in \cref{sec:intro}: an inference protocol
for dependent panels with measured size, an anchored relative null benchmarked on the real past,
oracle-based power comparisons that show where simpler tests win, a likelihood link and a
Markowitz study. The \WC{} test is one well-calibrated, interpretable tool within this protocol,
not a uniformly more powerful test.
